% Knowledge Graph Constructor: Evaluation (about 14 pages). Written in Stage 5.
% Every number has a row in evidence_registry.csv; sources are tests/gold/r62..r228 and REFACTOR_PLAN.md.
% Judged numbers: Claude Opus 5.5 (R66: Fable 5.1); rates are k/n with 95 % Wilson intervals where given.
\chapter{Knowledge Graph Constructor: Evaluation}\label{ch05}

This chapter evaluates the builder of \cref{ch04}. It answers two questions: whether the anchor graph is
\emph{correct as an index} --- its identity, links and provenance --- on four datasets (RQ1), and what the
measurements of the four representations say about storing facts against anchoring evidence (RQ2). The method is
stated first (\cref{sec:ch05-method}); the results follow by criterion, then by representation.

\section{Evaluation Questions and Method}\label{sec:ch05-method}

\paragraph{The graph is judged as an index.} The anchor graph's goal (\cref{sec:ch04-evolution}) is to locate
the things a need is about and lead from them to the evidence that concerns them. It is therefore not judged
by the answers of a question-answering system, which would mix the graph's quality with that of the reader, but
by ten criteria fixed in the design document before the first measurement (\cref{tab:criteria}). Navigation in
the measurements uses only the fixed walks W1--W5 over an offline snapshot of each build; no model takes part.
Where meaning must be judged --- whether a link is right, whether two merged names are one thing, whether a chunk
concerns its start --- the judge of \cref{sec:ch03-scoring-exact-match-and-judgement} decides, and code computes
the rates.

\begin{table}[htbp]
\centering
\caption{The builder's audit criteria. Hard criteria are gates: a failure is a defect. Comparative criteria are
reported and compared between designs.}
\label{tab:criteria}
\scriptsize
\renewcommand{\arraystretch}{1.25}
\begin{tabularx}{\linewidth}{@{}l>{\raggedright}p{2.2cm}>{\raggedright}X>{\raggedright\arraybackslash}p{3.4cm}@{}}
\toprule
& Criterion & Question & Rule \\
\midrule
C0 & Fidelity & Is the measured snapshot really the build? & \textbf{hard}: every logged count equal \\
C1 & Provenance & Does every item made from text have a chunk and quotation that show it? & \textbf{hard}: 1.0 \\
C2 & Findability & Does a target's name lead to its node (W1)? & comparative: hit@1, hit@5 \\
C3 & Identity & Are distinct things kept apart and one thing kept as one node? & \textbf{hard}: no wrong merge of records or individuals; concepts and splits comparative \\
C4 & Record linking & Does a text mention link to the right record? & \textbf{hard}: judged precision $\geq 0.95$, no confirmed link into another thing's scope \\
C5 & Evidence reach & From the right start, are the gold evidence chunks reached within $k$ chunks? & comparative: against vector retrieval at the same $k$ \\
C6 & Purity & Of the chunks reached from a record or individual, how many concern it? & \textbf{hard}: $\geq 0.95$; concept starts comparative \\
C7 & Selectivity & How much of the corpus does one start lead to? & reported: median and 90th percentile; a start reaching $>20\,\%$ is a hub \\
C8 & Connectivity & Are things the data connects reachable, with every hop witnessed? & \textbf{hard}: no hop without a record relation or a chunk \\
C9 & Generality, cost & Same code and settings everywhere; at what cost and size? & \textbf{hard}: no per-dataset code or setting \\
\bottomrule
\end{tabularx}
\end{table}

\paragraph{Builds.} The final graphs are those the query engine is evaluated on: furniture and generality as
rebuilt in step R108, NHTSA in R109 (a fix to derived claims that left every identity and link item unchanged),
and NTSB as built once, without any tuning, in R155b. The three earlier datasets were built with two extraction
passes, NTSB with one (\cref{sec:ch04-extraction}).

\paragraph{Gold.} The audit uses four reference sets, all written by Claude before the measurement:
\begin{itemize}[nosep]
\item for C2, C5 and C8, per question of the development sets, the names it starts from, the records it should reach
and its evidence chunks;
\item for C3, pairs of names that do or do not denote one thing;
\item for extraction, the claims a careful reader lists in 40 randomly chosen sentences per dataset;
\item for mentions, the mentions a reader marks in the same sentences.
\end{itemize}
The C3, C4 and C6 items are judged one by one: every merge, every mention-to-record link, and every pair of a start
and a chunk it reaches. The targets of C2, C5 and C8 were written after earlier builds had been inspected, so they
are not blind; this is stated wherever they are used.

\section{Extraction Quality}\label{sec:ch05-extraction}

The claims are no longer the graph's backbone, but they are still extracted and stored, and their quality is the
evidence for the design decision. Two measures are reported: \emph{precision} --- of the claims stored, how many
are right --- and \emph{recall} --- of the claims a reader finds in random sentences, how many are stored.

\paragraph{Precision is high.} \Cref{fig:claim-precision} shows that the content of a stored claim (its subject,
relation and object) is right in 96--99\,\% of cases on all four datasets: 289/293 on NTSB, 518/530 on furniture,
520/530 on NHTSA and 208/216 on generality. Strict precision, which also requires every assertion field to be
right (truth, modality, condition, time), is lower: 285/293, 498/530, 480/530 and 202/216. On the NHTSA data, 38
of the 59 recorded field faults are in the modality field --- a recall notice that a part \emph{may} fail stored as
a failure that happened.

\begin{sidewaysfigure}
\centering
\includegraphics[width=\linewidth]{results/construction/r229_build_accuracy/r229_fig1_claim_precision_by_dataset.pdf}
\caption{Precision of the stored extracted claims, judged against their chunks: strict (every field) and content
(the triple), with 95\,\% Wilson intervals. $n$ is every extracted claim of the build.}
\label{fig:claim-precision}
\end{sidewaysfigure}

\paragraph{Recall is low, mostly by design of the schema.} \Cref{fig:claim-recall} shows the other side. Of the
claims a reader lists in 40 random sentences, the builds store 42/72 on furniture (58\,\%), 19/67 on NHTSA, 12/66 on
generality and 13/74 on NTSB (18--28\,\%). The judge also recorded why each claim was missed. On NHTSA and
generality most misses had \emph{no fact type} that could hold them (32 of 48 and 39 of 54 misses): the schema the
model proposed for the goal did not cover statements the text makes about repairs, schedules or people. On NTSB
the schema could hold most of what was missed, and the model simply did not extract it (45 of 61 misses) --- the
build that used one extraction pass instead of two.

\begin{sidewaysfigure}
\centering
\includegraphics[width=\linewidth]{results/construction/r229_build_accuracy/r229_fig2_claim_recall_by_dataset.pdf}
\caption{A reader's claims in 40 random sentences per dataset: stored, missed although a fact type could hold
them, and missed for lack of a fact type.}
\label{fig:claim-recall}
\end{sidewaysfigure}

\paragraph{What this means for the design.} A graph whose stored facts are almost always right but cover only a
fifth to a half of what the text states cannot answer a question that counts or collects over the text: the
answer would be precise and incomplete, with nothing to signal the incompleteness. This is the quantitative core
of the decision in \cref{sec:ch04-evolution}, measured again on the final builds: a builder that must work on
any domain with a proposed schema should not make the stored facts the basis of exact answers.

\section{The Mention Index}\label{sec:ch05-mentions}

The mention index replaces the claims as the way into the text: every thing a question might start from should be
a mention. \Cref{tab:mentions} gives its recall against a reader's mentions in the same 40 sentences, and the
precision of a random sample of the mentions the pass found.

\begin{table}[htbp]
\centering
\caption{The mention index. Recall: by exact name (code) and with near names the judge maps (judge). Precision: a
seeded sample of the mentions the pass stored, judged.}
\label{tab:mentions}
\small
\begin{tabular}{@{}lcccc@{}}
\toprule
& NTSB & Furniture & NHTSA & Generality \\
\midrule
Recall, exact name & \kn{133}{184} & \kn{53}{75} & \kn{100}{125} & \kn{88}{109} \\
Recall, judged & \kn{153}{184} & \kn{68}{75} & \kn{116}{125} & \kn{103}{109} \\
Pass precision, judged & \kn{51}{57} & \kn{54}{59} & \kn{52}{59} & \kn{49}{60} \\
\bottomrule
\end{tabular}
\end{table}

Judged recall is 83--95\,\%, against 36--52\,\% for the names in the claims of the earlier build: the pass finds
the starting points the claims missed. Its precision, between 82 and 92\,\%, misses the bound of 0.90 set when the
pass was introduced on generality (49/60, interval 0.70--0.89) and, on a second sample, on NHTSA (52/59, interval
0.78--0.94; the first sample of the same build gave 55/59). The extra mentions are mostly generic words and
fragments (``night shift'', the identifier of a web page); since a mention only leads to chunks, an extra mention
costs reading, not a wrong fact.

\section{Identity and Record Linking}\label{sec:ch05-identity}

\paragraph{Final results.} On all four datasets every hard identity criterion passes (\cref{fig:audit}):
\begin{itemize}[nosep]
\item C4, links from mentions to records: 36/36 on NTSB, 63/63 on furniture, 77/77 on NHTSA and 48/48 on
generality, each above the 0.95 bound (lower Wilson bounds 0.90--0.95);
\item C3: no wrong merge of two records or two individuals on any dataset;
\item concept merges judged one thing: 79/80, 123/124, 66/68 and 52/53. Together with the individual merges,
this gives the C3 rates of the figure (80/81, 150/151, 93/95, 85/86). The NTSB rate leaves out two merges the
judge found ambiguous; counting them as wrong gives 80/83.
\end{itemize}

\begin{sidewaysfigure}
\centering
\includegraphics[width=\linewidth]{results/construction/r229_build_accuracy/r229_fig5_audit_by_dataset.pdf}
\caption{The judged audit criteria C3 (merges judged one thing), C4 (mention-to-record links right) and C6 (chunks
reached from a record or individual that concern it), with the 95\,\% bound of C4 and C6.}
\label{fig:audit}
\end{sidewaysfigure}

\paragraph{How the gates were reached.} These results were not the builder's first. \Cref{tab:trajectory} traces
the identity criteria from the first judged audit to the final builds. The first audit failed C4 on furniture: 7
of 49 links went to another product's part --- the desk's ``drawer'' linked to the nightstand's drawer assembly ---
and 7 records had been merged wrongly. Each step removed one cause, under a rule fixed beforehand: keep record
links inside the scope of the document's subject (R94); let code link only exact names (R95a); hand near misses to a
model whose choice code verifies (R95b, R96); judge individual identity on evidence (R100); turn stated kinds that
match no record into concepts (R107); and send a key made of plain words to the chooser rather than let it decide
alone (R108). The last step cost four correct NHTSA links (\cref{sec:ch05-limitations}): precision was chosen over
recall.

\begin{table}[htbp]
\centering
\caption{The identity criteria from the first judged audit to the final builds. C4: judged right / links made.
``Wrong merges'': records merged wrongly (furniture / NHTSA / generality). ``='': unchanged.}
\label{tab:trajectory}
\scriptsize
\begin{tabularx}{\linewidth}{@{}l>{\raggedright}Xcccc@{}}
\toprule
Step & Change & C4 furniture & C4 NHTSA & C4 generality & Wrong merges \\
\midrule
R93 & first judged audit & \kn{42}{49} (fails) & \kn{75}{76} & \kn{40}{40} & 7 / 1 / 0 \\
R94 & record links stay in the document's scope & \kn{42}{45} (fails) & = & = & 3 / 1 / 0 \\
R95a & code links only exact names & \kn{38}{38} & = & = & 0 / 1 / 0 \\
R95b & a model chooses among near misses, code verifies & \kn{49}{49} & = & = & 0 / 1 / 0 \\
R99 & candidates outside a scope & \kn{50}{50} & \kn{76}{77} & \kn{43}{43} & 0 / 1 / 0 \\
R102 & rebuild with the mention pass & \kn{64}{64} & \kn{76}{77} & \kn{47}{47} & 0 / 1 / 0 \\
R103 & rebuild with re-extracted claims & \kn{63}{63} & \kn{81}{83} & \kn{50}{50} & 0 / 2 / 0 \\
R107 & unmatched stated kinds become concepts & \kn{63}{63} & \kn{81}{83} & \kn{48}{48} & 0 / 2 / 0 \\
R108 & a plain-word key goes to the chooser & \kn{63}{63} & \kn{77}{77} & \kn{48}{48} & 0 / 0 / 0 \\
\bottomrule
\end{tabularx}
\end{table}

The trajectory is development on these datasets, and the final C3 and C4 values of furniture, NHTSA and
generality are measured on the same data the fixes were made for. NTSB is the exception: it was built once with the
frozen code, and passed C4 (36/36) and C3 without any change.

\paragraph{Identity pairs.} Against the reference pairs of names that do or do not denote one thing, the resolver
joined no two things that the reference keeps apart on NHTSA and generality (NTSB had no scorable pair of
different things), and one pair on furniture: it
merged ``dimmer'' with ``dimmer switch'' and ``dimmer function'', which the reference keeps apart but the blind judge
accepted. Recall --- pairs that should be joined and were --- is lower: 11/14 on furniture, 4/7 on NHTSA, 13/13 on
generality. The resolver errs towards keeping things apart.

\section{Provenance, Purity, Connectivity and Findability}\label{sec:ch05-provenance}

\Cref{tab:audit-final} gives the remaining criteria on the final builds.

\begin{table}[htbp]
\centering
\caption{The remaining criteria on the final builds (arm A, the anchor graph).}
\label{tab:audit-final}
\small
\begin{tabularx}{\linewidth}{@{}Xcccc@{}}
\toprule
& NTSB & Furniture & NHTSA & Generality \\
\midrule
C0 fidelity & pass & pass & pass & pass \\
C1 provenance & 1.0 & 1.0 & 1.0 & 1.0 \\
C6 purity, record and individual starts & \kn{58}{60} & \kn{109}{109} & \kn{86}{86} & \kn{100}{100} \\
C6 purity, concept starts & \kn{58}{58} & \kn{51}{51} & \kn{69}{69} & \kn{31}{31} \\
C8 gold connections reached & \kn{7}{7} & \kn{14}{14} & \kn{16}{16} & \kn{12}{12} \\
C8 unwitnessed hops & 0 & 0 & 0 & 0 \\
C2 targets found at rank 1 & \kn{62}{66} & \kn{72}{86} & \kn{53}{65} & \kn{45}{62} \\
C2 targets found in the first 5 & \kn{64}{66} & \kn{80}{86} & \kn{60}{65} & \kn{54}{62} \\
C7 median share of corpus per start & 2.9\,\% & 1.4\,\% & 1.2\,\% & 3.1\,\% \\
C7 hubs ($> 20\,\%$ of chunks) & 3 concepts & none & none & 4 \\
\bottomrule
\end{tabularx}
\end{table}

\paragraph{Provenance and purity.} Every item made from text has its chunk and quotation (C1 = 1.0), and the chunks
reached from a record or an individual concern it in every case on three datasets. On NTSB, two of the 60 chunks
reached from one start do not: the mention ``flight instructor'' covers three different accidents of one document,
because a mention is defined per document (\cref{sec:ch05-limitations}). At 58/60 (lower Wilson bound 0.89) the
criterion still passes its 0.95 point estimate, but the case shows a limit of the identity model on documents that
describe several events.

\paragraph{Connectivity.} Every connection that the multi-hop questions require is reachable, and no hop is
unwitnessed: every step between two things is a record relation or a chunk that names both.

\paragraph{Findability and selectivity.} The name lookup W1 finds 73--94\,\% of the targets at rank 1 and 87--97\,\%
in the first five; the misses are mostly names written differently from any node (``J.~Pike'' for a staff
member recorded as Jonathan Pike). A typical start leads to 1--3\,\% of its corpus. The hubs are starts that
legitimately concern much of a small corpus --- on generality, a pumping station that 11 of 32 documents discuss,
and one staff member and two pumps; on NTSB, the concepts ``engine'', ``right wing'' and ``runway''.

\section{Evidence Reach against Vector Retrieval}\label{sec:ch05-reach}

Criterion C5 asks whether the graph, walking from the right start, reaches a question's gold evidence chunks
within a fixed budget of $k$ chunks, compared with plain vector retrieval of the $k$ chunks most similar to the
question. It is measured in two ways: from the gold start nodes (the graph alone) and end to end, from the
question's names through W1. \Cref{tab:reach} gives the questions whose evidence is complete.

\begin{table}[htbp]
\centering
\caption{C5: questions whose gold evidence chunks are all reached within $k$ chunks, anchor graph against vector
retrieval; McNemar ``only graph / only vector'', exact $p$.}
\label{tab:reach}
\small
\begin{tabular}{@{}lccc@{}}
\toprule
& Furniture ($n=31$) & NHTSA ($n=28$) & Generality ($n=38$) \\
\midrule
gold start, $k=5$ & 20 vs 21 (2/3, 1.0) & 24 vs 26 (2/4, 0.69) & 31 vs 34 (1/4, 0.38) \\
gold start, $k=10$ & 29 vs 26 (5/2, 0.45) & 26 vs 26 (2/2, 1.0) & 36 vs 37 (1/2, 1.0) \\
end to end, $k=5$ & 17 vs 21 (2/6, 0.29) & 22 vs 26 (2/6, 0.29) & 27 vs 34 (1/8, \textbf{0.039}) \\
end to end, $k=10$ & 23 vs 26 (4/7, 0.55) & 24 vs 26 (2/4, 0.69) & 34 vs 37 (1/4, 0.38) \\
\bottomrule
\end{tabular}
\end{table}

The result is negative for the graph as a \emph{ranker}. From the right start it reaches the evidence about as
often as vector retrieval does; end to end, at the tight budget of five chunks, it trails, and significantly so on
generality, where the start is often a name written differently from any node. Without a budget the graph reaches
nearly everything (61/61, 36/36 and 51/51 gold chunks from the gold starts), but a walk that reaches everything is
not a ranking. The value of the graph shown by this chapter is therefore not that it finds evidence faster than
vector search, but that what it links is right: identity, record links and purity, which vector search does not
provide at all. The query engine (\cref{ch06}) draws the consequence: it uses the graph to find and verify the
things a question is about, and ranked text search to find passages.

\section{Comparison of Graph Representations}\label{sec:ch05-representations}

\Cref{tab:representations} collects the measurements behind the design evolution of \cref{sec:ch04-evolution}.

\begin{table}[htbp]
\centering
\caption{The four representations, as measured. Judged by meaning unless stated. QA: questions answered right.}
\label{tab:representations}
\scriptsize
\renewcommand{\arraystretch}{1.25}
\begin{tabularx}{\linewidth}{@{}>{\raggedright}p{2.6cm}>{\raggedright}Xccc@{}}
\toprule
Representation (step) & Measure & Furniture & NHTSA & Generality \\
\midrule
Claims as edges (R62) & claim precision & \kn{170}{170} & \kn{135}{146} & --- \\
& claim recall, reference triples & \kn{89}{96} & \kn{56}{68} & --- \\
& path truth (R63) & \kn{64}{100} & \kn{42}{56} & --- \\
\midrule
Claims as nodes (R64--R68) & path truth & \kn{89}{89} & \kn{107}{107} & --- \\
& coverage of a reader's claims (R68) & \kn{35}{72} & \kn{21}{67} & --- \\
\midrule
Layered model (R73) & QA: graph / records + vector / vector & 34 / 34 / 13 of 68 & 60 / 50 / 25 of 68 & 17 / 25 / 27 of 41 \\
& graph vs records + vector (McNemar) & 8/8, $p=1.0$ & 13/3, $p=0.021$ & 0/8, $p=0.008$ \\
\midrule
Anchor graph (R108--R109) & C4 record links & \kn{63}{63} & \kn{77}{77} & \kn{48}{48} \\
& C6 purity (records, individuals) & \kn{109}{109} & \kn{86}{86} & \kn{100}{100} \\
& claim precision, content (R110) & \kn{518}{530} & \kn{520}{530} & \kn{208}{216} \\
& claim recall, random sentences (R111) & \kn{42}{72} & \kn{19}{67} & \kn{12}{66} \\
\bottomrule
\end{tabularx}
\end{table}

Three caveats apply. The steps used different builds, and the R68 furniture graph was built with DeepSeek rather
than Gemini. The QA comparison of R73 used the reading and query stage of its time, not the agent of \cref{ch06}.
And the reference triples of R62 are a different measure from the random-sentence recall of R111: the first counts
facts of the schema's types in fully labelled documents, the second every claim a reader finds, which is why the
same builder scores 0.93 on the first and 0.58 on the second on furniture.

\paragraph{Navigating with or without claims.} The anchor design kept the claim layer in the build so that the
decision could be tested. Arm A navigates the final builds by the walks only; arm B also walks through the claims
(from a thing to its claims, and from a claim to its chunk and to the things at its ends). On the final builds the
two arms reach the gold evidence on exactly the same questions: no discordant question on furniture, NHTSA or
generality, and on NTSB a single discordant question in each of two modes, one favouring each arm. But arm B pays for this with exactly the hops the witness rule forbids:
132 of its 1{,}894 hops on furniture, 2{,}080 of 24{,}091 on NHTSA and 117 of 3{,}360 on NTSB are unwitnessed, so arm
B fails the C8 hard criterion on three of four datasets, and it creates the only hub of the NHTSA graph (a vehicle
make reaching 28\,\% of the chunks). It is also larger: 85.7 against 38.8 edges per chunk on furniture, 91.7 against
30.4 on NHTSA. On this data, navigating through claims adds no evidence and adds wrong connections.

\section{Generalisation}\label{sec:ch05-generalisation}

The criterion C9 requires that the same code and settings build every dataset; it holds by construction (no
dataset-specific code or setting exists) and is enforced by the domain-word tests of \cref{sec:ch04-engineering}.
Whether the results generalise is a different question. Three of the four datasets were used during development,
so their final scores are partly the result of fixes made for them. The NTSB dataset was not: it was chosen,
downloaded and given its question set after the builder's code was frozen, and built once. One general
capability was added just before its build because its tables needed it: keys made of several columns (step
R155a), which changes nothing for data with single-column keys. No prompt, threshold or rule was tuned on it. It passed every hard gate
(C1, C3, C4, C6, C8), with content precision 289/293 and judged mention recall 153/184. Its weaknesses are
informative: the lowest claim recall (13/74), from the single extraction pass, and the one identity failure caused
by documents that describe several accidents. The fifth dataset, which failed C4 (\cref{sec:ch03-datasets}), shows
the same kind of limit from another side: a generic word linked to the wrong one of several lighthouses named in
one article.

The four datasets share a form: short documents, each mostly about one or a few things, with names written in
full. Generalisation to other forms --- long narratives, scientific reports, large tables --- is untested.

\section{Cost and Size}\label{sec:ch05-cost}

\Cref{tab:build-cost} gives the model cost of the final builds at list prices; embedding calls are not priced
(\cref{sec:ch03-reproducibility}).

\begin{table}[htbp]
\centering
\caption{Model cost of the final builds (USD) and the anchor graph's size per chunk. The plan and text schema of
furniture, NHTSA and generality were fixed from an earlier build and are not included.}
\label{tab:build-cost}
\small
\begin{tabular}{@{}lcccc@{}}
\toprule
& NTSB & Furniture & NHTSA & Generality \\
\midrule
Extraction & 0.210 & 0.425 & 0.607 & 0.223 \\
Mention pass & 0.104 & 0.115 & 0.154 & 0.060 \\
Resolution & 0.150 & 0.156 & 0.656 & 0.081 \\
Plan and text schema & 0.176 & --- & --- & --- \\
\textbf{Total} & \textbf{0.640} & \textbf{0.696} & \textbf{1.417} & \textbf{0.365} \\
\midrule
Nodes / edges per chunk & 47.3 / 60.9 & 26.2 / 38.8 & 22.0 / 30.4 & 28.8 / 33.9 \\
\bottomrule
\end{tabular}
\end{table}

A build costs well under two dollars. The largest single item, NHTSA's resolution, is a known inefficiency: about
668 pairs of recall numbers that look alike in spelling were each put to the model and always answered ``different'',
at about \$1.1 in the first rebuild (\cref{sec:ch05-limitations}).

\section{Failure Analysis and Limitations}\label{sec:ch05-limitations}

\paragraph{Recorded limitations of the final builder.} When the refinement of the builder was closed, four
limitations were recorded rather than fixed:
\begin{enumerate}[nosep]
\item four correct NHTSA record links were lost by the last fix (R108): the chooser now refuses ``Civic Type R''
(twice) and two model-year ranges of the Civic, apparently reading the record's model year against the range ---
a loss of link recall, not precision;
\item two correct generality links (``the pump had run dry'') became concepts under the kind rule of R107;
\item identifiers are nominated as identity pairs by spelling --- a cost, not an error;
\item the ``dimmer'' ambiguity of \cref{sec:ch05-identity}.
\end{enumerate}

\paragraph{Open issues found later.} Mention identity is per document. In the NTSB reports, each of which
describes five accidents, ``the pilot'' or ``flight instructor'' of different accidents becomes one mention; this
causes the two C6 failures and makes 24 of the 27 NTSB reference pairs unscorable. A fix --- per-section identity for
documents that describe several events --- was designed but not built, as the builder was frozen. Other open items
are minor: generic words among the pass's mentions, and near-duplicate mentions (``may not engage'', ``not engage'').

\paragraph{Limitations of the evaluation.} The reference sets and the judge come from one model family, with no
agreement figure against a human judge. The targets of C2, C5 and C8 are not blind to earlier builds. Three of four
datasets were used for development. Each corpus is small. These limits are discussed with those of the query
engine in \cref{sec:ch08-threats-to-validity}.

\section{Discussion}\label{sec:ch05-discussion}

\paragraph{RQ1: can the pipeline build a correct, provenance-aware, entity-centric graph without domain-specific
code?} On the four datasets measured, yes, within stated limits. The same code built all four graphs. Every item
made from text has its verified quotation; no two records or individuals were merged wrongly; every judged link from
a mention to a record was right; and walks from a record or an individual reached only chunks that concern it, with
one exception on documents describing several events. These gates were reached through development on three of
the datasets and held without changes on the fourth, unseen one. The fifth dataset that failed C4 shows that they
are not guaranteed.

\paragraph{RQ2: what are the trade-offs between storing facts and anchoring evidence?} The measurements point one
way. Stored facts are precise (96--99\,\% content precision) but incomplete (18--58\,\% recall against a reader),
mostly because a proposed schema cannot anticipate what a corpus says; facts stored on shared nodes leaked across
entities (36 of 100 paths false); and a rich layered model of facts did not improve answers over records plus
vector retrieval. The anchor graph gives up exact answers from stored facts --- which the recall figures show were
not available anyway --- and in exchange makes identity, record links and provenance exact and checkable.
Navigating through claims added no evidence and added unwitnessed connections. The graph is not a better ranker of
evidence than vector search (C5); its contribution is to say \emph{what} a passage is about and \emph{which} record
a name refers to, correctly. How a reader can use that index to answer questions is the subject of
\cref{ch06,ch07}.
